% §1.6 Notes and Lineage
%
% Cross-references: sec:three-faces, sec:quantifying-noise, sec:sampling-rates,
% sec:why-atom, sec:bsc, sec:miller-rabin, sec:bib-notes-ch3,
% ch:classification-measures. DEFINED: sec:bib-notes-ch1.
% Citations: shannon-1948, bross-1954, warner-1965,
%   erlingsson-pihur-korolova-2014, manning-raghavan-schutze,
%   towell-bernoulli-sets, towell-bernoulli-data-type.

\section{Notes and Lineage}
\label{sec:bib-notes-ch1}

\paragraph{Lineage.}
Every ingredient of the noisy bit is classical. The channel reading, a
conditional distribution of output given input, and the binary symmetric
channel are Shannon's \parencite{shannon-1948}. The $2 \times 2$ table of sensitivity
and specificity comes from diagnostic testing, where \textcite{bross-1954}
studied how misclassification distorts it. The same asymmetric channel applied
deliberately to a survey respondent's true indicator is randomized
response \parencite{warner-1965}, deployed at scale in RAPPOR
\parencite{erlingsson-pihur-korolova-2014}; there the channel is known and
inverted to estimate a prevalence, the reverse of the error-accounting use made
of it here. Precision and recall come from information retrieval
\parencite[Ch.~8]{manning-raghavan-schutze}. The binomial estimators and
intervals of \Cref{sec:sampling-rates} are standard and are developed, with
their sources, in chapter~7; Miller-Rabin is credited in the
\hyperref[sec:bib-notes-ch3]{notes to chapter~3}. The framing of these objects as the atom of a
framework of approximate types is that of the source manuscripts
\parencite{towell-bernoulli-sets,towell-bernoulli-data-type}.

\paragraph{What this chapter adds.}
A presentation: the biased coin, the noisy line, and the classifier verdict are
treated as three faces of one atom with two parameters, and the distinction
between the intrinsic pair $(\fprate, \fnrate)$ and the workload-dependent pair
(precision, recall) is fixed as the book's convention. No result in the chapter
is new.

\paragraph{Left out.}
ROC analysis and threshold selection, cost-sensitive decision rules, and the
inversion of a known channel to estimate prevalence (the randomized-response
problem) are mentioned only in passing.
